
% BEGIN COMPLETE CD AND DEGREE STEP 20260921

\clearpage\begingroup
\part*{Every boundary-kernel entry and its exact compressed correction}
\newcommand{\cdrootC}{\mathbb C}
\newcommand{\cdrootR}{\mathbb R}
\newcommand{\cdrootconj}[1]{\overline{#1}}


\section{Source, mass, coordinates and scope}

This is a finite calculation on the original Gamma receiver MR22--23
\cite{cdroot:MR}%
% BEGIN VERIFIED PUBLIC PROOF LINK
~\href{https://github.com/KokunoYumeto/zeta-function-research-reader/blob/5992ad50c0f2a06921f1fcaded7b4e38097c0739/workbenches/splitzero-tandem/continuations/20260920-original-kernel-mixed-determinants/ORIGINAL_KERNEL_MIXED_DETERMINANT.tex\#L36}{[proof MR1]}%
% END VERIFIED PUBLIC PROOF LINK
{}, not a replacement of the arithmetic source or its original
coefficient frame. The underlying Christoffel--Darboux identity is
classical; Akemann--Vernizzi's original author source explicitly gives
its real-line formula and its algebraic continuation in the footnote
at \texttt{CD} \cite{cdroot:AV}. The exact Gamma recurrence and norms below
are the Meixner--Pollaczek formulas \cite{cdroot:DLMF}, transported in the
original JM26 and KG6 coordinates \cite{cdroot:JM,cdroot:KG}%
% BEGIN VERIFIED PUBLIC PROOF LINK
~\href{https://github.com/KokunoYumeto/zeta-function-research-reader/blob/5992ad50c0f2a06921f1fcaded7b4e38097c0739/workbenches/splitzero-tandem/continuations/20260920-original-kernel-mixed-determinants/providers/JOINT_MINIMUM_RECEIVERS.tex\#L286}{[proof JM21]} \href{https://github.com/KokunoYumeto/zeta-function-research-reader/blob/5992ad50c0f2a06921f1fcaded7b4e38097c0739/workbenches/splitzero-tandem/continuations/20260920-original-kernel-mixed-determinants/providers/KERNEL_GAMMA_INDEPENDENT.tex\#L73}{[proof KG1]}%
% END VERIFIED PUBLIC PROOF LINK
{}. We derive every finite
identity used here. No large-degree coefficient, limit estimate,
arithmetic zero statement, or priority claim about the classical
identity is asserted.

The retained density, original physical coordinate and its complex
evaluation chart are
\begin{equation}\label{cdroot:eq:source}
 d\sigma(y)=\frac{|\Gamma(1/4+iy/2)|^2}{2\pi}\,dy,
 \quad M_\sigma=\sqrt{2\pi},\quad c=k/2,
 \quad S=c+iy,
 \quad u=\frac{v-c}{i},\quad w=\frac{z-c}{i}.
 \tag{CD1}
\end{equation}
In formulas with two scalar polynomial variables, write \(x=u\) and
\(t=\overline w\); these variables are independent during
differentiation. An asterisk on a matrix always means conjugate
transpose. The source inner product is
\(\langle P,Q\rangle=\int\overline{P(c+iy)}Q(c+iy)\,d\sigma(y)\).

Here is the mass check, retaining the scale in \(y\). The original DLMF
weight and norm at \(\lambda=1/4,\phi=\pi/2\) are
\[
 |\Gamma(1/4+ix)|^2\,dx,
 \qquad
 h_j=\frac{2\pi\Gamma(j+1/2)}{\sqrt2\,j!}.
\]
Set \(y=2x\) and
\(p_j(y)=j!P_j^{(1/4)}(y/2;\pi/2)\).
The measure becomes \(|\Gamma(1/4+ix)|^2dx/\pi\). Hence the actual
squared norm is
\begin{equation}\label{cdroot:eq:norm}
 \gamma_j=\frac{(j!)^2}{\pi}h_j
 =\sqrt2\,j!\Gamma(j+1/2)
 =\sqrt{2\pi}\,j!(1/2)_j.
 \tag{CD2}
\end{equation}
In particular \(\gamma_0=M_\sigma=\sqrt{2\pi}\). The exact density
change \(dy=2dx\) is already included in CD2.

The original generating function, after the same substitution, is
\[
 \sum_{j\geq0}\frac{p_j(y)}{j!}\zeta^j
 =(1+\zeta^2)^{-1/4}e^{y\arctan\zeta}.
\]
Its highest \(y\) coefficient proves monicity. Differentiating in
\(\zeta\) gives
\((1+\zeta^2)\partial_\zeta F=(y-\zeta/2)F\).
Coefficient comparison therefore proves the original recurrence
\begin{equation}\label{cdroot:eq:recurrence}
 p_0=1,\quad p_1=y,\quad
 p_{j+1}(y)=yp_j(y)-a_jp_{j-1}(y),
 \quad a_j=j(j-\tfrac12)=\frac{\gamma_j}{\gamma_{j-1}}
 \quad(j\geq1).
 \tag{CD3}
\end{equation}
All coefficients of \(p_j\) are real. We put \(p_{-1}=0\) only in
recurrences whose coefficient is zero; no negative-index norm is used.

\section{The two-polynomial identity and its sign}

For \(d\geq0\), set
\[
 \mathcal K_d(x,t)=\sum_{j=0}^{d-1}\frac{p_j(x)p_j(t)}{\gamma_j},
 \qquad K_d(v,\overline z)=\mathcal K_d(u,\overline w).
\]
An empty sum is zero. For \(d\geq1\), define
\[
 F_d(x,t)=p_d(x)p_{d-1}(t)-p_{d-1}(x)p_d(t).
\]
Then the exact polynomial identity is
\begin{equation}\label{cdroot:eq:cd}
 (x-t)\mathcal K_d(x,t)=\frac{F_d(x,t)}{\gamma_{d-1}}.
 \tag{CD4}
\end{equation}
To prove it without dividing by \(x-t\), put
\(T_j=p_{j+1}(x)p_j(t)-p_j(x)p_{j+1}(t)\).
The recurrence gives
\[
 (x-t)p_j(x)p_j(t)=T_j-a_jT_{j-1}.
\]
For \(j=0\) the second term is zero. Division by \(\gamma_j\) and
\(a_j/\gamma_j=1/\gamma_{j-1}\) make the finite sum telescope to
\(T_{d-1}/\gamma_{d-1}=F_d/\gamma_{d-1}\).
This proves \eqref{cdroot:eq:cd} as an identity in \(\cdrootC[x,t]\), including
\(x=t\). Away from that diagonal it gives
\begin{equation}\label{cdroot:eq:offdiagonal}
 \boxed{\mathcal K_d(x,t)=
 \frac{p_d(x)p_{d-1}(t)-p_{d-1}(x)p_d(t)}
      {\gamma_{d-1}(x-t)}.}
 \tag{CD5}
\end{equation}
At the diagonal, differentiation of the polynomial identity in \(x\)
and then setting \(x=t\) gives
\begin{equation}\label{cdroot:eq:diagonal}
 \boxed{\mathcal K_d(t,t)=
 \frac{p_d'(t)p_{d-1}(t)-p_{d-1}'(t)p_d(t)}{\gamma_{d-1}}.}
 \tag{CD6}
\end{equation}
At \(d=1\) both formulas give \(1/M_\sigma\), fixing the denominator
sign. At \(d=0\), the kernel and all its derivatives are zero; neither
CD5 nor CD6 is used.

The collision in these formulas is \(u=\overline w\), not merely
\(u=w\). For complex off-line arguments, the Hermitian kernel entry
\(\mathcal K_d(u,\overline u)\) usually belongs to CD5, not CD6.
CD6 applies when the two polynomial arguments coincide, whether or not
that common argument is real. Positivity is asserted for Hermitian Gram
matrices, not for every analytically continued scalar entry.

\section{Every divided jet, both away from and at a collision}

For a polynomial \(f\), write \(f^{[a]}=f^{(a)}/a!\). Define
\[
 F_{d;a,b}(x,t)=p_d^{[a]}(x)p_{d-1}^{[b]}(t)
             -p_{d-1}^{[a]}(x)p_d^{[b]}(t),
 \qquad
 \mathcal K_{d;r,s}=\frac{\partial_x^r\partial_t^s\mathcal K_d}{r!s!}.
\]
For all nonnegative integers \(r,s\), at \(x\ne t\),
\begin{equation}\label{cdroot:eq:jets-away}
 \boxed{\mathcal K_{d;r,s}(x,t)=\frac1{\gamma_{d-1}}
  \sum_{a=0}^r\sum_{b=0}^s
  \frac{(-1)^a\binom{a+b}{a}\,
        F_{d;r-a,s-b}(x,t)}{(x-t)^{a+b+1}}.}
 \tag{CD7}
\end{equation}
Indeed the divided \((a,b)\) derivative of \((x-t)^{-1}\) is
\((-1)^a\binom{a+b}{a}(x-t)^{-a-b-1}\). The divided-derivative product
rule is convolution with no further binomial factor. Applying it to
CD5 proves CD7.

At a collision \((x,t)=(\tau,\tau)\), the exact formula is
\begin{equation}\label{cdroot:eq:jets-collision}
 \boxed{\mathcal K_{d;r,s}(\tau,\tau)
  =\frac1{\gamma_{d-1}}
      \sum_{j=0}^{s}F_{d;r+1+j,s-j}(\tau,\tau)
  =-\frac1{\gamma_{d-1}}
      \sum_{j=0}^{r}F_{d;r-j,s+1+j}(\tau,\tau).}
 \tag{CD8}
\end{equation}
For a complete proof, expand both sides of
\(F_d(x,t)=\gamma_{d-1}(x-t)\mathcal K_d(x,t)\) at
\(x=\tau+h,t=\tau+z\). Equality of the coefficient of \(h^az^b\)
gives
\[
 F_{d;a,b}(\tau,\tau)=\gamma_{d-1}
 \bigl(\mathcal K_{d;a-1,b}(\tau,\tau)
       -\mathcal K_{d;a,b-1}(\tau,\tau)\bigr),
\]
where any negative derivative index means zero. Starting with
\((a,b)=(r+1,s)\) and iterating downward in \(s\) proves the first
sum. Starting with \((r,s+1)\) and iterating downward in \(r\)
proves the second. The expansion is finite, so there is no interchange
of infinite limiting operations. CD7--8 include arbitrary repeated
roots and all mixed divided jets. Derivatives beyond the degree are
zero, exactly as in the original finite kernel.

\section{The original \(S\) coordinate and its reflection denominator}

The monic polynomials in the physical coordinate are
\[
 P_j(v)=i^jp_j((v-c)/i),\qquad \|P_j\|^2=\gamma_j.
\]
The paired phases cancel in each kernel summand, so this is the same
\(K_d\) already defined. Since
\[
 u-\overline w=-i(v+\overline z-2c),
\]
and
\[
 F_d(u,\overline w)=-i\bigl(
 P_d(v)\overline{P_{d-1}(z)}+
 P_{d-1}(v)\overline{P_d(z)}\bigr),
\]
CD5 becomes
\begin{equation}\label{cdroot:eq:S-CD}
 \boxed{K_d(v,\overline z)=
 \frac{P_d(v)\overline{P_{d-1}(z)}+
       P_{d-1}(v)\overline{P_d(z)}}
      {\gamma_{d-1}(v+\overline z-2c)}.}
 \tag{CD9}
\end{equation}
This expression is used away from \(v+\overline z=2c\); the
polynomial extension supplies the remaining values. The numerator is a
sum and the denominator is a reflection sum. Replacing either by the
usual real-line difference would change the original coordinate.

The exact divided-jet transport is
\begin{equation}\label{cdroot:eq:S-jets}
 \frac{\partial_v^r\partial_{\overline z}^sK_d(v,\overline z)}{r!s!}
   =i^{-r}i^s\mathcal K_{d;r,s}(u,\overline w).
 \tag{CD10}
\end{equation}
This follows from \(\partial_v=i^{-1}\partial_u\) and
\(\partial_{\overline z}=i\partial_{\overline w}\), so it also
fixes the phases at collisions. For a direct formula in physical
variables, write \(G_d(v,b)=P_d(v)\overline{P_{d-1}(\overline b)}+
P_{d-1}(v)\overline{P_d(\overline b)}\), a polynomial in independent
variables \(v,b\), and use divided derivatives \(G_{d;a,b}\), with
the second subscript here an integer. Away from \(v+b=2c\), the
counterpart of CD7 has denominator \((v+b-2c)^{a+b+1}\) and sign
\((-1)^{a+b}\) in each summand. At a point \(v_0+b_0=2c\), it is
\begin{equation}\label{cdroot:eq:S-jets-collision}
 \frac{\partial_v^r\partial_b^sK_d(v,b)}{r!s!}
 \bigg|_{(v_0,b_0)}
  =\frac1{\gamma_{d-1}}
      \sum_{j=0}^s(-1)^jG_{d;r+1+j,s-j}(v_0,b_0).
 \tag{CD11}
\end{equation}
To prove this independently, expand
\(G_d=\gamma_{d-1}(v+b-2c)K_d\) at that point; the local linear
factor is now \(h+z\), so solving the same coefficient recurrence
alternates the signs. This agrees with CD10.

\section{Every entry at the original quartet and all four cutoffs}

The original simple-packet domain is
\[
 0<\delta<\tfrac12,\quad \gamma>2,\quad k\equiv1\pmod4,
 \quad k\geq9,\quad q=(k+1)^2,
\]
on the original five-orbit period locus of OCF \cite{cdroot:OCF}%
% BEGIN VERIFIED PUBLIC PROOF LINK
~\href{https://github.com/KokunoYumeto/zeta-function-research-reader/blob/5992ad50c0f2a06921f1fcaded7b4e38097c0739/workbenches/splitzero-tandem/continuations/20260920-original-kernel-mixed-determinants/providers/ORIGINAL_CONDUCTOR_STRIP_FRAME.tex\#L27}{[proof OCF1]}%
% END VERIFIED PUBLIC PROOF LINK
{}. Here
\(\gamma\) is the retained zero height, whereas \(\gamma_j\) is a
squared polynomial norm. The ordered roots and their exact charts are
\begin{equation}\label{cdroot:eq:roots}
 v_{ab}=c+(2a-k)\delta+i(2b-k)\gamma,
 \qquad u_{ab}=(2b-k)\gamma-i(2a-k)\delta,
 \quad 0\leq a,b\leq k.
 \tag{CD12}
\end{equation}
Distinctness follows by comparing the two real coordinates, since
\(\delta,\gamma>0\). However
\[
 u_{ab}=\overline{u_{a'b'}}
 \quad\Longleftrightarrow\quad b=b',\quad a'=k-a.
\]
Indeed equality of real parts gives \(b=b'\), and equality of imaginary
parts gives \(-(2a-k)=2a'-k\). Thus the full covariance has exactly
one collision entry per row, at its reflected column \((k-a,b)\).
There are \(q\) such ordered entries. Since \(k\) is odd, there are no
fixed points \(a=k-a\), and these are not its Hermitian diagonal
entries. Their CD6 values must be retained.

For \(d=N+1\), let \(C_d\) denote the matrix
\([\mathcal K_d(u_\alpha,\overline{u_\beta})]\) in exactly that
lexicographic root order. For every pair of entries it is computed from
the same two polynomials by
\begin{equation}\label{cdroot:eq:entry-recipe}
 (C_d)_{\alpha\beta}=
 \begin{cases}
 \dfrac{p_d(u_\alpha)\overline{p_{d-1}(u_\beta)}-
 p_{d-1}(u_\alpha)\overline{p_d(u_\beta)}}
 {\gamma_{d-1}(u_\alpha-\overline{u_\beta})},
 &u_\alpha\ne\overline{u_\beta},\\[7pt]
 \dfrac{p_d'(u_\alpha)p_{d-1}(u_\alpha)-
 p_{d-1}'(u_\alpha)p_d(u_\alpha)}{\gamma_{d-1}},
 &u_\alpha=\overline{u_\beta}.
 \end{cases}
 \tag{CD13}
\end{equation}
This is an exact two-boundary-polynomial expression, including the
derivatives of those same polynomials at the resonant entries.
The four original cutoffs require exactly these pairs:
\[
\begin{array}{c|c|c|c}
N&d&\text{boundary pair}&\text{denominator norm}\\\hline
q-1&q&(p_q,p_{q-1})&\gamma_{q-1}\\
q&q+1&(p_{q+1},p_q)&\gamma_q\\
2q-1&2q&(p_{2q},p_{2q-1})&\gamma_{2q-1}\\
2q&2q+1&(p_{2q+1},p_{2q})&\gamma_{2q}
\end{array}
 \tag{CD14}
\]
Every number in the table follows from \(d=N+1\); no degree shift is
absorbed into a different kernel convention.

The original nonreal frame from OCF and MR6--8 is retained literally:
\[
 Z=[\,M_A\ \ S_kZ_k\,]\in\cdrootC^{q\times r},\qquad
 \pi Z=0,\quad r=k^2-6k+17,
 \quad m=q-r=8k-16.
\]
All its period and amplitude coefficients, pivot choices and column
order remain those of the original construction. The compressed matrix
is therefore given entry by entry by
\begin{equation}\label{cdroot:eq:compression-recipe}
 A_d=Z^TC_d\overline Z,
 \qquad
 (A_d)_{ij}=\sum_{\alpha,\beta=1}^q
    Z_{\alpha i}(C_d)_{\alpha\beta}\overline{Z_{\beta j}},
 \tag{CD15}
\end{equation}
with CD13 substituted in every term. Thus two boundary polynomials
suffice for the full original compression at each cutoff, provided all
root entries, all nonreal frame coefficients and all collision values
are retained. This statement is not an assertion that the two already
compressed boundary vectors alone obey a closed rank-two displacement.
The exact additional term is calculated next.

\section{Full displacement and its confluent version}

Put \(Y=\operatorname{diag}(u_\alpha)\) and
\(f_j=(p_j(u_\alpha))_\alpha\). CD4 is precisely the matrix identity
\begin{equation}\label{cdroot:eq:full-displacement}
 YC_d-C_dY^*=
 \frac{f_df_{d-1}^*-f_{d-1}f_d^*}{\gamma_{d-1}}
 =:B_d.
 \tag{CD16}
\end{equation}
The right side is skew-Hermitian and has rank at most two. For a
resonant entry the scalar multiplier on the left is zero, and the
corresponding numerator on the right is zero. Consequently CD16 alone
does not determine those entries of \(C_d\); CD6 or CD8 supplies them.
For example a matrix supported at a single resonant entry belongs to
the kernel of \(X\mapsto YX-XY^*\). This proves the need for the
collision data without assuming uniqueness for a singular matrix map.

For arbitrary finite divided-jet blocks, define instead
\[
 f_j=\bigl(p_j^{[a]}(u_\alpha)\bigr)_{(\alpha,a)},
 \qquad J=\bigoplus_\alpha(u_\alpha I+L_\alpha),
 \qquad (L_\alpha)_{a,a-1}=1.
\]
The other entries of each \(L_\alpha\) are zero; rows are ordered by
increasing derivative degree. Divided Leibniz differentiation gives
\((yp_j)^{[a]}(u_\alpha)=u_\alpha p_j^{[a]}(u_\alpha)
+p_j^{[a-1]}(u_\alpha)\). Hence
\(Jf_j=f_{j+1}+a_jf_{j-1}\), and the same telescoping proof gives
\begin{equation}\label{cdroot:eq:jet-displacement}
 JC_d-C_dJ^*=
 \frac{f_df_{d-1}^*-f_{d-1}f_d^*}{\gamma_{d-1}},
 \quad C_d=\sum_{j<d}\frac{f_jf_j^*}{\gamma_j}.
 \tag{CD17}
\end{equation}
Its entries are exactly the divided derivatives in CD7--8. If the total
jet count is at most \(d\), the first that many monic polynomial columns
have full row rank: a polynomial of smaller degree in the kernel would
be divisible by the product of all the specified root powers, by its
finite Taylor expansions, and would therefore be zero. Thus \(C_d>0\)
in the original admitted-degree regime. The displacement identity itself
does not require this positivity or a matrix inverse.

\section{Exact compression and the complementary coupling term}

For the original simple-root frame put
\[
\begin{gathered}
 V=\overline Z,\quad W=V^*V=Z^T\overline Z,\quad P=VW^{-1}V^*,\quad Q=I-P,\\
 H=(V^*YV)W^{-1},\quad g_j=V^*f_j=Z^Tf_j.
\end{gathered}
\]
Full column rank of \(Z\), proved in OCF and MR7, makes \(W>0\).
Direct multiplication gives \(P^*=P\), \(P^2=P\),
\(PV=V\), and \(V^*Q=0\). Therefore \(P\) is the Euclidean
orthogonal projection to \(\operatorname{im}V\). This projection is
an algebraic device; the original source Gram and the original columns
of \(Z\) are not changed.

The placement of \(W^{-1}\) is essential: \(A_d=V^*C_dV\) is a
covariant Gram, so the reduced operator appearing on its left is
\(H=(V^*YV)W^{-1}\). The coefficient operator with the opposite
placement \(W^{-1}(V^*YV)\) cannot be inserted into the following
formula without also transporting the Gram.

Define the two actual complement maps
\[
 L=QY^*V,\qquad T_d^{\perp}=QC_dV,\qquad
 \mathcal E_d=L^*T_d^{\perp}-(T_d^{\perp})^*L.
\]
Then the exact compressed displacement is
\begin{equation}\label{cdroot:eq:compressed-displacement}
 \boxed{HA_d-A_dH^*=
 \frac{g_dg_{d-1}^*-g_{d-1}g_d^*}{\gamma_{d-1}}
   -\mathcal E_d.}
 \tag{CD18}
\end{equation}
To prove it, use
\[
 HA_d=V^*YPC_dV,\qquad A_dH^*=V^*C_dPY^*V.
\]
Insert \(P=I-Q\), apply CD16 to the full terms, and retain the
remaining difference
\(V^*YQC_dV-V^*C_dQY^*V=L^*T_d^\perp-(T_d^\perp)^*L\).
This proves both the sign and every factor in CD18.

Equivalently, put \(\rho_j=V^*YQf_j\). Projecting the original
three-term recurrence gives its exact inhomogeneous receiver
\[
 Hg_j=g_{j+1}+a_jg_{j-1}-\rho_j.
\]
Since \(T_d^\perp=\sum_{j<d}Qf_jg_j^*/\gamma_j\), its correction is
\begin{equation}\label{cdroot:eq:correction-sum}
 \mathcal E_d=\sum_{j<d}
    \frac{\rho_jg_j^*-g_j\rho_j^*}{\gamma_j}.
 \tag{CD19}
\end{equation}
Thus all lower-degree complement couplings occur, with their exact
signs, in addition to the two projected boundary vectors. Formula CD15
still computes this entire matrix from the two full boundary
polynomials and their collision data; CD19 describes what a closed
recurrence on compressed vectors alone would otherwise omit.

The original complementary frame is available without selecting a new
one. Put
\[
 U=\pi^T\in\cdrootC^{q\times m},\qquad M=U^*U=\overline\pi\pi^T.
\]
The coefficient identity \(\pi Z=0\) gives \(U^*V=0\); the ranks
sum to \(q\), so \(Q=UM^{-1}U^*\). Define
\[
 F=V^*YU,\qquad T_d=U^*C_dV.
\]
Substitution into the already proved correction gives
\begin{equation}\label{cdroot:eq:actual-complement}
 \boxed{\mathcal E_d=FM^{-1}T_d-T_d^*M^{-1}F^*,
 \qquad
 \operatorname{rank}(HA_d-A_dH^*)\leq\min(r,2+2m).}
 \tag{CD20}
\end{equation}
Each correction summand has rank at most \(m\), so their difference
has rank at most \(2m\). Adding the rank-at-most-two boundary term
proves the bound. Equivalently the displacement's image lies in the
span of the \(2+2m\) columns
\([g_d,g_{d-1},F,T_d^*]\). The exact original dimensions are
\(r=k^2-6k+17\), \(m=8k-16\); consequently the finite bound is
\(\min(k^2-6k+17,16k-30)\). This is a finite rank statement, not an
estimate of a determinant or a large-degree return.
More explicitly, in block sizes \(1,1,m,m\), put
\(\mathcal G_d=[g_d,g_{d-1},F,T_d^*]\). Then the same formula is
\[
 HA_d-A_dH^*=\mathcal G_d
 \begin{pmatrix}
 0&\gamma_{d-1}^{-1}&0&0\\
 -\gamma_{d-1}^{-1}&0&0&0\\
 0&0&0&-M^{-1}\\
 0&0&M^{-1}&0
 \end{pmatrix}\mathcal G_d^*.
\]
Expanding its four nonzero blocks gives CD18 and CD20 exactly. The
matrix \(M\) depends only on the original coefficient frame, not on
the degree cutoff or on an inverse of the complete source covariance.

For a fixed cutoff the rank-two formula without the correction holds
if and only if \(L^*T_d^\perp\) is Hermitian, equivalently
\(FM^{-1}T_d=T_d^*M^{-1}F^*\). The stronger equation \(L=0\) is
sufficient at every cutoff and is equivalent to
\(Y^*\operatorname{im}V\subseteq\operatorname{im}V\).
For the original diagonal matrix \(Y\), whose eigenvalues are distinct,
this is equivalent to invariance under \(Y\): polynomial interpolation
at the finite distinct eigenvalues expresses \(Y^*\) as a polynomial
in \(Y\), and conversely expresses \(Y\) as a polynomial in \(Y^*\).
Neither invariance is asserted for the original OCF image merely from
its definition as an invariant-polynomial coefficient image. CD18--20
apply to its exact coefficients whether or not the displayed correction
vanishes.

All compression calculations also apply to CD17 with \(Y\) replaced
by its divided-jet multiplication matrix \(J\). For that possibly
nonnormal matrix, \(L=0\) means \(J^*\)-invariance. The diagonal
interpolation argument does not apply to a Jordan block, and no
equivalence with \(J\)-invariance is used.
In this general divided-jet statement \(r\) denotes the retained frame's
column count and \(m\) its complementary dimension. The OCF polynomial
rank formulas above are claimed only on the original simple-quartet
domain, not newly extended to a repeated-root coefficient construction.

\section{A finite Gamma example in which the correction is nonzero}

The following is an exact algebra fixture, not a sample of the original
period. It demonstrates why a general compression argument cannot
delete the term in CD18. Take the actual Gamma recurrence coefficient
\(a_1=1/2\), let \(M=M_\sigma\), and use
\[
 u=(0,i),\quad Y=\operatorname{diag}(0,i),\quad
 V=(1,i)^T,\quad Z=(1,-i)^T,\quad d=2.
\]
Then \(\gamma_0=M\), \(\gamma_1=M/2\),
\(f_0=(1,1)^T\), \(f_1=(0,i)^T\), and
\(f_2=(-1/2,-3/2)^T\). Direct multiplication gives
\[
 C_2=M^{-1}\begin{pmatrix}1&1\\1&3\end{pmatrix},\qquad
 A_2=4/M,\qquad W=2,\qquad H=i/2.
\]
Therefore \(HA_2-A_2H^*=4i/M\). But
\(g_1=1\), \(g_2=-1/2+3i/2\), so the compressed boundary term is
\[
 (g_2\overline{g_1}-g_1\overline{g_2})/\gamma_1=6i/M.
\]
The difference is the nonzero correction \(\mathcal E_2=2i/M\), with
the minus sign in CD18. This example has nonreal frame entries and
uses the retained positive Gamma mass. It establishes failure of an
uncorrected compression formula, not a claim about the value of the
original period's correction.

\section{Physical-coordinate displacement and mass powers}

Let \(X=cI+iY=\operatorname{diag}(v_\alpha)\),
\(h_j=i^jf_j\), and \(H_S=cI+iH\).
Multiplying CD16 by \(i\) gives the exact physical-coordinate identity
\begin{equation}\label{cdroot:eq:S-displacement}
 XC_d+C_dX^*-2cC_d
 =\frac{h_dh_{d-1}^*+h_{d-1}h_d^*}{\gamma_{d-1}}.
 \tag{CD21}
\end{equation}
Indeed \(h_dh_{d-1}^*=if_df_{d-1}^*\) and
\(h_{d-1}h_d^*=-if_{d-1}f_d^*\). The corresponding original-frame
compression is
\begin{equation}\label{cdroot:eq:S-compressed}
 H_SA_d+A_dH_S^*-2cA_d
 =\frac{(i^dg_d)(i^{d-1}g_{d-1})^*
       +(i^{d-1}g_{d-1})(i^dg_d)^*}{\gamma_{d-1}}
       -i\mathcal E_d.
 \tag{CD22}
\end{equation}
Every term is Hermitian, since \(\mathcal E_d\) is skew-Hermitian.
Thus both the anti-Hermitian displacement in the Gamma coordinate and
the shifted Hermitian displacement in the physical coordinate are
retained exactly.

As a diagnostic only, replacing \(d\sigma\) by \(a\,d\sigma\),
\(a>0\), leaves all monic polynomials, root coordinates, frames and
projection matrices unchanged. It multiplies every \(\gamma_j\) by
\(a\), so \(C_d,A_d,T_d,T_d^\perp,\mathcal E_d\) all scale by
\(a^{-1}\). Both sides of every CD formula scale identically.
The original source uses \(a=1\) and the exact mass in CD1--2.

\section{Verification and source boundary}

The companion checker compares finite polynomial sums against CD5--11,
the full and divided-jet displacement identities, the actual matrix
form of CD18--20 in declared nonreal algebra fixtures, the collision
pattern of the original quartet coordinate formula, the four cutoff
shifts, and the retained mass factors. Those fixtures do not evaluate
the actual period or change the original frame in the theorem.
The separate mathematical audit retains a second exact compression
derivation and its explicit nonzero-correction check.

\begin{thebibliography}{9}
\bibitem{cdroot:AV} G. Akemann and G. Vernizzi,
\emph{Characteristic Polynomials of Complex Random Matrix Models},
arXiv:hep-th/0212051v2, original author source
\texttt{AkemannVernizzi.tex}, \texttt{CD} and its following footnote,
\url{https://arxiv.org/abs/hep-th/0212051v2}.
\bibitem{cdroot:DLMF} T. H. Koornwinder, R. Wong, R. Koekoek and R. F. Swarttouw;
W. P. Reinhardt, \emph{Orthogonal Polynomials}, NIST DLMF Chapter 18,
original retained TeX for equations
\url{https://dlmf.nist.gov/18.19.E8},
\url{https://dlmf.nist.gov/18.19.E9a}, and
\url{https://dlmf.nist.gov/18.23.E7}.
The calculation above retains the exact density change \(dy=2dx\).
\bibitem{cdroot:MR} Split-Zero programme,
\emph{The original observation kernel as an exact mixed-relation
determinant}, 20 September 2026, complete original MR1--23.

% BEGIN VERIFIED PUBLIC PROOF LINK
\href{https://github.com/KokunoYumeto/zeta-function-research-reader/blob/5992ad50c0f2a06921f1fcaded7b4e38097c0739/workbenches/splitzero-tandem/continuations/20260920-original-kernel-mixed-determinants/ORIGINAL_KERNEL_MIXED_DETERMINANT.tex\#L36}{Source proof MR1}.
% END VERIFIED PUBLIC PROOF LINK
\bibitem{cdroot:OCF} Split-Zero programme,
\emph{An exact conductor frame and a fixed invariant-generator
recursion}, 14 September 2026, complete original OCF1--26,
especially OCF17 and OCF21--24.

% BEGIN VERIFIED PUBLIC PROOF LINK
\href{https://github.com/KokunoYumeto/zeta-function-research-reader/blob/5992ad50c0f2a06921f1fcaded7b4e38097c0739/workbenches/splitzero-tandem/continuations/20260920-original-kernel-mixed-determinants/providers/ORIGINAL_CONDUCTOR_STRIP_FRAME.tex\#L27}{Source proof OCF1}.
% END VERIFIED PUBLIC PROOF LINK
\bibitem{cdroot:JM} Split-Zero programme,
\emph{The joint source in the original arithmetic metric: full
spectrum, observation, and heat return}, original JM21--28,
especially JM26--27.

% BEGIN VERIFIED PUBLIC PROOF LINK
\href{https://github.com/KokunoYumeto/zeta-function-research-reader/blob/5992ad50c0f2a06921f1fcaded7b4e38097c0739/workbenches/splitzero-tandem/continuations/20260920-original-kernel-mixed-determinants/providers/JOINT_MINIMUM_RECEIVERS.tex\#L286}{Source proof JM21}.
% END VERIFIED PUBLIC PROOF LINK
\bibitem{cdroot:KG} Split-Zero programme,
\emph{Independent derivation of the complete Gamma reduction for the
original kernel}, 20 September 2026, KG1--21; and the accompanying
\emph{COERCIVITY\_INDEPENDENT}, C1--4, retaining the original DLMF
TeX sources and their norm conversion.

% BEGIN VERIFIED PUBLIC PROOF LINK
\href{https://github.com/KokunoYumeto/zeta-function-research-reader/blob/5992ad50c0f2a06921f1fcaded7b4e38097c0739/workbenches/splitzero-tandem/continuations/20260920-original-kernel-mixed-determinants/providers/KERNEL_GAMMA_INDEPENDENT.tex\#L73}{Source proof KG1}.
% END VERIFIED PUBLIC PROOF LINK
\end{thebibliography}

\endgroup

\clearpage\begingroup
\part*{Independent complete compression, jet and metric calculation}
\providecommand{\cdaudittightlist}{\setlength{\itemsep}{0pt}\setlength{\parskip}{0pt}}

\section{Exact compressed Christoffel--Darboux displacement}\label{cdaudit:exact-compressed-christoffeldarboux-displacement}

Independent algebra audit, 2026-09-20. All results are finite-dimensional. The original roots, frame, and Gram matrix remain fixed. No assertion about generic or actual OCF period values is made. The scalar Gamma recurrence and full CD identity are mathematical inputs supplied by the parent derivation; the compression statements below are proved here. The degree-two fixture uses the supplied exact Gamma coefficients and checks its full identity directly.

\subsection{CD-C1. Fixed objects and the full boundary}\label{cdaudit:cd-c1.-fixed-objects-and-the-full-boundary}

Let \(1\le r\le N\), \(d\ge1\), and retain \[
u_\alpha=(v_\alpha-c)/i,\qquad Y=\operatorname{diag}(u_\alpha),\qquad
Z\in\mathbb C^{N\times r},\quad \operatorname{rank}Z=r.
\] Put \[
V=\overline Z,\quad W=V^*V=Z^T\overline Z,\quad
C=\sum_{j=0}^{d-1}\frac{f_jf_j^*}{\gamma_j},\quad
A=V^*CV=Z^TC\overline Z,
\] where \(f_j=(p_j(u_\alpha))_\alpha\), \(\gamma_j>0\), \(^{*}\) means conjugate transpose, and \(T\) means transpose without conjugation. Then \(C=C^*\ge0\), \(A=A^*\ge0\), and \[
YC-CY^*=\frac{f_df_{d-1}^*-f_{d-1}f_d^*}{\gamma_{d-1}}.
\tag{CD-C1}
\] The full boundary has rank at most two.

The finite telescoping proof of the supplied boundary is worth recording. For monic polynomials with the actual recurrence \[
xp_j=p_{j+1}+\alpha_jp_j+\beta_jp_{j-1},\quad
\alpha_j\in\mathbb R,\quad \beta_j=\gamma_j/\gamma_{j-1}\ (j\ge1),
\] and \(p_{-1}=0\), the \(j\)-th contribution to \(YC-CY^*\) is \[
\frac{f_{j+1}f_j^*-f_jf_{j+1}^*}{\gamma_j}
+\frac{f_{j-1}f_j^*-f_jf_{j-1}^*}{\gamma_{j-1}}.
\] The second term is absent at \(j=0\); the real \(\alpha_j\)-terms cancel. Summation cancels all internal pairs and leaves (CD-C1). This does not replace any coefficient of the original Gamma recurrence.

The Gram \(W\) is positive definite because \(x^*Wx=\|Vx\|^2>0\) for nonzero \(x\). Define \[
P=VW^{-1}V^*,\quad Q=I_N-P,\quad S=\operatorname{ran}V.
\] Direct multiplication gives \(P^*=P\), \(P^2=P\), and \(PV=V\). Thus \(P\) is the Euclidean orthogonal projection onto \(S\), and \(Q\) projects onto \(S^\perp\).

\subsection{CD-C2. Exact compression and the correct operator orientation}\label{cdaudit:cd-c2.-exact-compression-and-the-correct-operator-orientation}

Define matrices and maps \[
B=V^*YV,\quad H=BW^{-1},\quad
L=QY^*V:\mathbb C^r\to S^\perp,\quad
G=QCV:\mathbb C^r\to S^\perp,
\] \[
E=L^*G=V^*YQCV,\quad
\Delta=E-E^*=L^*G-G^*L,\quad g_j=V^*f_j=Z^Tf_j.
\] Then the exact result in the original frame is \[
\boxed{HA-AH^*
=\frac{g_dg_{d-1}^*-g_{d-1}g_d^*}{\gamma_{d-1}}-\Delta.}
\tag{CD-C2}
\] Proof: \(HA=V^*YPCV\), while \(AH^*=V^*CPY^*V\). Inserting \(I=P+Q\) gives \[
\begin{aligned}
V^*(YC-CY^*)V
&=V^*YPCV-V^*CPY^*V\\
&\quad+V^*YQCV-V^*CQY^*V\\
&=HA-AH^*+\Delta.
\end{aligned}
\] Compression of (CD-C1) supplies the stated boundary. The identity \(E=L^*G\) uses \(Q^*=Q=Q^2\).

The primal coordinate matrix for \(PY|_S\) is \(K=W^{-1}B\), since \(PYV=VK\). The left operator in (CD-C2) is instead \(H=BW^{-1}=WKW^{-1}\). Equivalently, for \[
K_\sharp=W^{-1}V^*Y^*V,
\] one has \(PY^*V=VK_\sharp\), \(H=K_\sharp^*\), and the left side is \(K_\sharp^*A-AK_\sharp\). No invariant-subspace assertion was used.

The placement can also be checked without changing the original frame. Under a hypothetical \(V'=VR\), \(R\) invertible, \[
W'=R^*WR,\ A'=R^*AR,\ B'=R^*BR,\quad
K'=R^{-1}KR,\quad H'=R^*HR^{-*}.
\] Thus \[
H'A'-A'H'^*=R^*(HA-AH^*)R,\qquad \Delta'=R^*\Delta R.
\] This proves congruence covariance for the Gram identity.

\subsection{CD-C3. Exact vanishing and nonvanishing criteria}\label{cdaudit:cd-c3.-exact-vanishing-and-nonvanishing-criteria}

For the given \(C\), the correction-free formula holds if and only if \[
\boxed{L^*G=G^*L.}
\tag{CD-C3}
\] This follows directly from (CD-C2). Equivalently, \[
\langle Lx,Gy\rangle=\langle Gx,Ly\rangle
\quad\text{for all }x,y\in\mathbb C^r,
\] or \[
\operatorname{Im}\langle Lx,Gx\rangle=0
\quad\text{for all }x\in\mathbb C^r.
\tag{CD-C4}
\] Here \(\langle a,b\rangle=a^*b\). To prove the second equivalence, note \(x^*\Delta x=2i\operatorname{Im}(x^*Ex)\). If every such quadratic form is zero, applying it to \(e_j\), \(e_j+e_k\), and \(e_j+ie_k\) forces each matrix entry of \(\Delta\) to vanish. Conversely \(\Delta=0\) annihilates those quadratic forms. In particular, one \(x\) with nonzero imaginary part in (CD-C4) proves nonclosure.

Exact sufficient mechanisms are:

\begin{itemize}
\cdaudittightlist
\item
  \(L=0\) if and only if \(Y^*S\subset S\), since \(Q\) annihilates exactly \(S\).
\item
  \(G=0\) if and only if \(CS\subset S\).
\item
  Orthogonality of \(\operatorname{ran}L\) and \(\operatorname{ran}G\) gives \(L^*G=0\).
\end{itemize}

Each makes \(\Delta=0\). Neither leakage need vanish when their pairing is Hermitian. Thus nonzero leakage alone does not prove failure.

For diagonal \(Y\), \(YS\subset S\) and \(Y^*S\subset S\) are equivalent. Indeed, let \(\lambda_1,\ldots,\lambda_t\) be its distinct diagonal values, and define \[
h(z)=\sum_{a=1}^t\overline{\lambda_a}
\prod_{b\ne a}\frac{z-\lambda_b}{\lambda_a-\lambda_b}.
\] Then \(Y^*=h(Y)\). Invariance under \(Y\) implies invariance under every power and hence under \(h(Y)\). Repeating with \(Y^*\) gives the converse. This proof does not establish either invariance for the actual OCF subspace.

For clarity about the scope of the obstruction, the following universal algebra statement is also exact: \(\Delta=0\) for every positive-definite Hermitian ambient \(C\) if and only if \(L=0\). Only necessity remains to prove. If \(L\ne0\), set \(G_0=iL\) and \[
C_0=G_0W^{-1}V^*+VW^{-1}G_0^*.
\] Then \(C_0=C_0^*\), \(G_0^*V=0\), and \(C_0V=G_0\). Choose real \(\lambda>\|C_0\|_{\rm op}\), and set \(C=C_0+\lambda I>0\); positivity follows from \[
z^*Cz\ge(\lambda-\|C_0\|_{\rm op})\|z\|^2>0.
\] Since \(QV=0\), \(QCV=G_0\), so \(\Delta=2iL^*L\ne0\). This statement concerns general Hermitian matrices; it asserts no genericity of the actual Gamma kernels or OCF periods.

\subsection{CD-C4. Receiver complement and precise rank bound}\label{cdaudit:cd-c4.-receiver-complement-and-precise-rank-bound}

Let \(\pi\in\mathbb C^{m\times N}\) have full row rank, \(\pi Z=0\), and \(N=r+m\), and put \(U=\pi^T\). In the parent\textquotesingle s receiver, \(m=8k-16\). Then \[
U^*V=\overline\pi\,\overline Z=\overline{\pi Z}=0.
\] Since \(\operatorname{rank}U=m=\dim S^\perp\), \(\operatorname{ran}U=S^\perp\). Define \[
M_U=U^*U>0,\quad F=V^*YU\in\mathbb C^{r\times m},
\quad T=U^*CV\in\mathbb C^{m\times r}.
\] The subscript distinguishes the receiver Gram \(M_U\) from scalar mass \(M\). The projection proof from CD-C1 gives \[
Q=UM_U^{-1}U^*,\quad E=FM_U^{-1}T,\quad
\boxed{\Delta=FM_U^{-1}T-T^*M_U^{-1}F^*.}
\tag{CD-C5}
\] In particular, the signs in the compressed displacement are boundary minus the first product plus its adjoint.

There is an explicit generator factorization \[
HA-AH^*=\mathcal B\mathcal J\mathcal B^*,\qquad
\mathcal B=[\,g_d\ \ g_{d-1}\ \ F\ \ T^*\,],
\] \[
\mathcal J=
\begin{pmatrix}
0&\gamma_{d-1}^{-1}&0&0\\
-\gamma_{d-1}^{-1}&0&0&0\\
0&0&0&-M_U^{-1}\\
0&0&M_U^{-1}&0
\end{pmatrix},
\tag{CD-C6}
\] with block sizes \(1,1,m,m\). Multiplication gives exactly (CD-C2) and (CD-C5). Since the image is contained in that of the \(r\)-by-\((2+2m)\) matrix \(\mathcal B\), \[
\boxed{\operatorname{rank}(HA-AH^*)\le\min(r,2+2m),\quad m=8k-16.}
\tag{CD-C7}
\] Likewise \(\operatorname{rank}\Delta\le2m\). These are upper bounds, not equality assertions. If \(m=0\), the receiver blocks are empty, \(Q=0\), and the bound is \(\min(r,2)\).

\subsection{CD-C5. Nonreal exact Gamma degree-two algebra fixture}\label{cdaudit:cd-c5.-nonreal-exact-gamma-degree-two-algebra-fixture}

Keep arbitrary scalar mass \(M>0\). The original supplied value is \(M=\sqrt{2\pi}\), not \(\sqrt2\,\pi\). Use \[
p_0=1,\quad p_1=x,\quad p_2=x^2-\tfrac12,\qquad
\gamma_0=M,\quad \gamma_1=M/2,
\] and \[
u=(0,i),\quad v=(c,c-1),\quad Y=\operatorname{diag}(0,i),
\quad Z=\binom1{-i},\quad V=\overline Z=\binom1i.
\] The coordinate identity \(u=(v-c)/i\) holds exactly. This is a declared algebra fixture, not an evaluation of the actual period-derived OCF frame.

The vectors and complete kernel are \[
f_0=\binom11,\quad f_1=\binom0i,\quad f_2=\binom{-1/2}{-3/2},
\quad C=\frac1M\begin{pmatrix}1&1\\1&3\end{pmatrix}>0.
\] The leading principal minors are \(1/M\) and \(2/M^2\), proving positivity. Direct multiplication gives \[
D=YC-CY^*=\frac1M\begin{pmatrix}0&i\\i&6i\end{pmatrix}
=\frac{f_2f_1^*-f_1f_2^*}{M/2}.
\] The compressed data are \[
W=2,\quad A=4/M,\quad B=i,\quad H=i/2,
\] \[
V^*DV=6i/M,\qquad HA-AH^*=4i/M,\qquad
\boxed{\Delta=2i/M\ne0.}
\tag{CD-C8}
\] Explicitly, \[
P=\frac12\begin{pmatrix}1&-i\\i&1\end{pmatrix},\quad
L=\binom{i/2}{1/2},\quad G=\frac1M\binom{i-1}{1+i},
\quad E=(1+i)/M.
\] Hence \(E-E^*=2i/M\), exactly the missing term. A compatible receiver fixture is \[
\pi=(i,1),\ U=\binom i1,\ M_U=2,\ F=1,\ T=(2+2i)/M.
\] It satisfies \(\pi Z=0\), \(U^*V=0\), and \(FM_U^{-1}T=(1+i)/M\).

More generally \(p_2=x^2-\beta\), \(\gamma_1=M\beta\), \(\beta>0\), gives \[
A=(2+\beta^{-1})/M,\quad
V^*DV=2i(1+\beta^{-1})/M,\quad
HA-AH^*=i(2+\beta^{-1})/M,\quad
\Delta=i/(M\beta).
\] The Gamma fixture is its exact specialization \(\beta=1/2\).

\subsection{CD-C6. Literal all-jet displacement and factorial map}\label{cdaudit:cd-c6.-literal-all-jet-displacement-and-factorial-map}

Let \(\xi_\alpha\) be distinct nodes with multiplicities \(n_\alpha\), \(N=\sum n_\alpha\), ordered by \(\alpha\) and then \(k=0,\ldots,n_\alpha-1\). Define \[
\mathcal T_{\rm der}p=(p^{(k)}(\xi_\alpha))_{\alpha,k}.
\] Multiplication by \(x\) has the literal-jet matrix \[
J_{\rm der}=\bigoplus_\alpha J_{\alpha,\rm der},\quad
(J_{\alpha,\rm der})_{kk}=\xi_\alpha,\quad
(J_{\alpha,\rm der})_{k,k-1}=k\ (k\ge1),
\tag{CD-C9}
\] all other entries zero. This follows from \[
(xp)^{(k)}(\xi_\alpha)
=\xi_\alpha p^{(k)}(\xi_\alpha)+k p^{(k-1)}(\xi_\alpha),
\] with no second term at \(k=0\). Consequently \(\mathcal T_{\rm der}(xp)=J_{\rm der}\mathcal T_{\rm der}p\). Writing \(f_{j,\rm der}=\mathcal T_{\rm der}p_j\) and \[
C_{\rm der}=\sum_{j<d}f_{j,\rm der}f_{j,\rm der}^*/\gamma_j,
\] the same fully finite recurrence cancellation proves \[
J_{\rm der}C_{\rm der}-C_{\rm der}J_{\rm der}^*
=\frac{f_{d,\rm der}f_{d-1,\rm der}^*
-f_{d-1,\rm der}f_{d,\rm der}^*}{\gamma_{d-1}}.
\tag{CD-C10}
\] All identities CD-C2--CD-C7 follow with these objects and the original literal-coordinate frame, because their proofs used no normality.

Let \(D_{\rm fac}=\operatorname{diag}(k!)_{\alpha,k}\). The exact divided-derivative map is \[
\mathcal T_{\rm div}=D_{\rm fac}^{-1}\mathcal T_{\rm der},\quad
J_{\rm div}=D_{\rm fac}^{-1}J_{\rm der}D_{\rm fac}.
\] Its subdiagonal entries are \(1\). Moreover \[
f_{j,\rm der}=D_{\rm fac}f_{j,\rm div},\qquad
C_{\rm der}=D_{\rm fac}C_{\rm div}D_{\rm fac}.
\tag{CD-C11}
\] The same compressed form requires the dual frame relation \[
V_{\rm div}=D_{\rm fac}V_{\rm der},\quad
V_{\rm der}^*C_{\rm der}V_{\rm der}
=V_{\rm div}^*C_{\rm div}V_{\rm div}.
\] This is an exact relation between presentations, not a replacement of the fixed frame. Their Euclidean Gram matrices need not agree, so original \(W,P,H\) must either be computed in the original coordinates or transported with the metric below.

\subsection{CD-C7. Jordan invariance is not adjoint invariance}\label{cdaudit:cd-c7.-jordan-invariance-is-not-adjoint-invariance}

For jets, \(L=0\) is exactly \(J_{\rm der}^*S\subset S\). Ordinary \(J_{\rm der}\)-invariance alone is insufficient. A full all-jet Gamma fixture proves this.

At the single node \(\xi=i\), with derivative orders \(0,1\), use the same degree-two polynomials and norms as CD-C5: \[
J=\begin{pmatrix}i&0\\1&i\end{pmatrix},\
f_0=\binom10,\ f_1=\binom i1,\ f_2=\binom{-3/2}{2i},\
C=\frac1M\begin{pmatrix}3&2i\\-2i&2\end{pmatrix}>0.
\] Its positive leading minors are \(3/M\) and \(2/M^2\). Direct multiplication, or CD-C10, verifies the full boundary. Choose \(V=e_2\). Then \(S=\mathbb Ce_2\) is \(J\)-invariant, but \[
W=1,\ A=2/M,\ H=i,\quad
L=\binom10,\quad G=\frac1M\binom{2i}0.
\] Therefore \[
\Delta=4i/M\ne0,\quad
V^*(JC-CJ^*)V=8i/M,\quad HA-AH^*=4i/M.
\tag{CD-C12}
\] This fixture proves the stated insufficiency; it asserts nothing about actual confluent OCF period values.

\subsection{CD-C8. Polynomial quotient identity and exact metric transport}\label{cdaudit:cd-c8.-polynomial-quotient-identity-and-exact-metric-transport}

Let \(q(x)=\prod_\alpha(x-\xi_\alpha)^{n_\alpha}\). Division by the monic degree-\(N\) polynomial \(q\) gives a unique remainder of degree \(<N\): successively subtract its leading multiples for existence; uniqueness follows because a nonzero multiple of \(q\) cannot have degree \(<N\). Use the original ordered coefficient basis \(1,x,\ldots,x^{N-1}\) for these remainders. Let \(R_q\) be multiplication by \(x\) modulo \(q\), and define the literal-jet matrix \[
\mathsf T_{(\alpha,k),\ell}
=\begin{cases}
\ell!\,\xi_\alpha^{\ell-k}/(\ell-k)!,&\ell\ge k,\\
0,&\ell<k,
\end{cases}\qquad 0\le\ell<N.
\] It is invertible: a degree-\(<N\) polynomial with zero indicated jets is divisible by every \((x-\xi_\alpha)^{n_\alpha}\), and hence their product, so is zero. Both spaces have dimension \(N\), so this injective map is bijective. Jets agree for polynomials differing by a multiple of \(q\), by the product rule. The multiplication formula therefore proves \[
\mathsf T R_q=J_{\rm der}\mathsf T.
\tag{CD-C13}
\]

Let \(a_j\) be the coefficient vector of the remainder of \(p_j\), and put \[
K_q=\sum_{j<d}a_ja_j^*/\gamma_j.
\] Then \(f_{j,\rm der}=\mathsf T a_j\) and \(C_{\rm der}=\mathsf T K_q\mathsf T^*\). Multiplying CD-C10 by \(\mathsf T^{-1}\) and \(\mathsf T^{-*}\) proves \[
R_qK_q-K_qR_q^*
=\frac{a_da_{d-1}^*-a_{d-1}a_d^*}{\gamma_{d-1}}.
\tag{CD-C14}
\]

Retain the original literal-jet frame \(V\). Set \[
\widehat V=\mathsf T^*V,\quad
\mathsf G=\mathsf T^{-1}\mathsf T^{-*}=(\mathsf T^*\mathsf T)^{-1}.
\] Since \(V=\mathsf T^{-*}\widehat V\), multiplication gives \[
A=\widehat V^*K_q\widehat V,\quad
W=\widehat V^*\mathsf G\widehat V,\quad
B=\widehat V^*R_q\mathsf G\widehat V.
\tag{CD-C15}
\] The original operator and correction are consequently \[
H=(\widehat V^*R_q\mathsf G\widehat V)
(\widehat V^*\mathsf G\widehat V)^{-1},
\] \[
E=\widehat V^*R_q
\bigl(I-\mathsf G\widehat VW^{-1}\widehat V^*\bigr)
K_q\widehat V,\qquad \Delta=E-E^*.
\tag{CD-C16}
\] To verify the latter formula, use \[
\mathsf T^{-1}Q\mathsf T
=I-\mathsf G\widehat VW^{-1}\widehat V^*
\] in \(E=V^*J_{\rm der}QC_{\rm der}V\), followed by CD-C13. These are exact maps between the polynomial and all-jet presentations. Replacing the \(W\) of CD-C15 by \(\widehat V^*\widehat V\) would change the original metric.

\subsection{Verification and limits}\label{cdaudit:verification-and-limits}

The accompanying exact checker tests the fixtures, nonscalar nonorthonormal orientation, receiver factorization, literal and divided jets, polynomial quotient transport, and negative controls. The proofs above establish the identities independently of those checks.

For the fixed original OCF \(Z\), the unconditional conclusions are CD-C2, the necessary-and-sufficient criterion CD-C3, the exact receiver CD-C5, and rank bound CD-C7. This audit does not claim the actual period correction vanishes or is nonzero. No frozen sibling file, source index, rendering, or publication was changed.

\endgroup

\clearpage\begingroup
\part*{Degree-step determinants and the original invariant-image cancellation}
\newcommand{\dsrootC}{\mathbb C}
\newcommand{\dsrootR}{\mathbb R}
\newcommand{\dsrootcalD}{\mathscr D}


\section{The retained objects and the finite outcome}

This note calculates one-degree updates of the original Gamma covariance,
its compression by the full original invariant-image frame, and the
inverse restriction to the original observation kernel. It uses the
finite Christoffel--Darboux calculation CD1--22 \cite{dsroot:CD}%
% BEGIN VERIFIED PUBLIC PROOF LINK
~\href{https://github.com/KokunoYumeto/zeta-function-research-reader/blob/main/workbenches/splitzero-tandem/continuations/20260921-kernel-degree-steps/COMPLETE_CD_BOUNDARY.tex\#L42}{[proof CD1]}%
% END VERIFIED PUBLIC PROOF LINK
{}, with its
classical source in Akemann--Vernizzi's original \texttt{CD} formula
\cite{dsroot:AV}, and the original Meixner--Pollaczek source identities
\cite{dsroot:DLMF}. It retains the original OCF coefficient construction
\cite{dsroot:OCF}%
% BEGIN VERIFIED PUBLIC PROOF LINK
~\href{https://github.com/KokunoYumeto/zeta-function-research-reader/blob/5992ad50c0f2a06921f1fcaded7b4e38097c0739/workbenches/splitzero-tandem/continuations/20260920-original-kernel-mixed-determinants/providers/ORIGINAL_CONDUCTOR_STRIP_FRAME.tex\#L27}{[proof OCF1]}%
% END VERIFIED PUBLIC PROOF LINK
{} and the original MR8--23 receiver \cite{dsroot:MR}%
% BEGIN VERIFIED PUBLIC PROOF LINK
~\href{https://github.com/KokunoYumeto/zeta-function-research-reader/blob/5992ad50c0f2a06921f1fcaded7b4e38097c0739/workbenches/splitzero-tandem/continuations/20260920-original-kernel-mixed-determinants/ORIGINAL_KERNEL_MIXED_DETERMINANT.tex\#L36}{[proof MR1]}%
% END VERIFIED PUBLIC PROOF LINK
{}.
Every determinant and inverse identity used below is proved directly.
No large-degree coefficient or estimate at the unresolved \(kq\) scale
is asserted.

Keep the actual simple-quartet, five-orbit domain
\[
\begin{gathered}
0<\delta<\tfrac12,\quad \gamma>2,\quad k\equiv1\pmod4,\quad k\geq9,\\
c=k/2,\quad q=(k+1)^2,\quad r=k^2-6k+17,\quad m=8k-16,\quad q=r+m.
\end{gathered}
\]
The zero height is \(\gamma\); the polynomial norms below are
\(\gamma_j\). In the original lexicographic order,
\begin{equation}\label{dsroot:eq:roots}
v_{ab}=c+(2a-k)\delta+i(2b-k)\gamma,
\quad u_{ab}=(v_{ab}-c)/i
=(2b-k)\gamma-i(2a-k)\delta,
\quad 0\leq a,b\leq k.
\tag{DS1}
\end{equation}
The original scalar relation is \(\chi(S)=\prod_{a,b}(S-v_{ab})\).
The Gamma comparison source, never divided by its mass, is
\begin{equation}\label{dsroot:eq:gamma}
\begin{gathered}
d\sigma(y)=\frac{|\Gamma(1/4+iy/2)|^2}{2\pi}\,dy,\quad M_\sigma=\sqrt{2\pi},\\
p_0=1,\quad p_1=y,\quad p_{j+1}=yp_j-j(j-\tfrac12)p_{j-1},\\
\gamma_j=\|p_j\|_\sigma^2=M_\sigma j!(1/2)_j.
\end{gathered}
\tag{DS2}
\end{equation}
These are monic orthogonal polynomials. To check the exact mass against
the retained DLMF formulas \cite[18.19.E8--E9a,18.23.E7]{dsroot:DLMF}, set
\(\lambda=1/4,\phi=\pi/2,x=y/2\), and multiply the degree-\(j\)
Meixner--Pollaczek polynomial by \(j!\). The norm
\(2\pi\Gamma(j+1/2)/(\sqrt2\,j!)\) in \(dx\), multiplied by
\((j!)^2/\pi\), is precisely \(\gamma_j\). The generating function
\((1+z^2)^{-1/4}\exp(y\arctan z)\) gives the displayed recurrence
by differentiating in \(z\). The physical monic polynomial is
\(P_j(S)=i^jp_j((S-c)/i)\), with the same squared norm.

Use the complete original frame and quotient, not the conductor alone:
\begin{equation}\label{dsroot:eq:frames}
\begin{gathered}
Z=[M_A\ \ S_kZ_k]\in\dsrootC^{q\times r},\quad \pi\in\dsrootC^{m\times q},\quad \ker\pi=\operatorname{im}Z,\\
V=\overline Z,\quad U=\pi^T,\quad W=V^*V,\quad M_U=U^*U.
\end{gathered}
\tag{DS3}
\end{equation}
OCF7, OCF17 and OCF21--24 prove these exact ranks and the coefficient
maps, with their original period, amplitude and pivot coefficients.
The dual coefficient pairing is bilinear, hence the transpose in
\(U=\pi^T\). In particular \(U^*V=\overline{\pi Z}=0\), and
\(U,V\) are complementary full-column-rank frames.

For \(d=N+1\geq q\), define
\begin{equation}\label{dsroot:eq:covariance}
f_j=(p_j(u_{ab}))_{ab},\quad g_j=V^*f_j=Z^Tf_j,
\quad C_d=\sum_{j<d}\frac{f_jf_j^*}{\gamma_j},
\quad A_d=V^*C_dV,\quad G_d=U^*C_d^{-1}U.
\tag{DS4}
\end{equation}
Here \(A_d\) is the compressed Gram asked for, and \(G_d\) is the
original restricted least-source-norm metric at cutoff \(N=d-1\).
The first \(q\) monic evaluation columns are a Vandermonde matrix
times a unit-triangular matrix; their determinant is nonzero because
the roots are distinct. Thus \(C_d,A_d,G_d\) are positive definite.
The exact minimum formula follows by solving the complete evaluation
constraint in the orthogonal source basis: if its coefficient matrix
is \(B\), then \(BB^*=C_d\), and \(B^*C_d^{-1}y\) is a lift of
\(y\) orthogonal to \(\ker B\). Every other lift adds an orthogonal
kernel vector, and hence increases its squared norm by that vector's
squared norm. Restriction by the original value-space frame \(U\)
therefore gives exactly \(G_d\), after minimizing all source relations.

\section{Entry updates, including every reflected entry}

The definition gives the entire one-degree update before any compression:
\begin{equation}\label{dsroot:eq:rankone}
C_{d+1}=C_d+\frac{f_df_d^*}{\gamma_d},\qquad
A_{d+1}=A_d+\frac{g_dg_d^*}{\gamma_d}.
\tag{DS5}
\end{equation}
No entry is omitted. For precision, the full CD evaluation is
\[
(C_d)_{\alpha\beta}=
\frac{p_d(u_\alpha)\overline{p_{d-1}(u_\beta)}
-p_{d-1}(u_\alpha)\overline{p_d(u_\beta)}}
{\gamma_{d-1}(u_\alpha-\overline{u_\beta})}
\]
away from \(u_\alpha=\overline{u_\beta}\). Its proof is the finite
telescoping of the recurrence in DS2, using
\(\gamma_j/\gamma_{j-1}=j(j-1/2)\). At a collision its exact value is
\(\mathcal W_d(u_\alpha)/\gamma_{d-1}\), where
\[
\mathcal W_d(x)=p_d'(x)p_{d-1}(x)-p_{d-1}'(x)p_d(x).
\]
In the original quartet, equality of the two real coordinates in DS1
proves
\(u_{ab}=\overline{u_{a'b'}}\) exactly when
\((a',b')=(k-a,b)\). There is one such entry in each of the \(q\)
rows, and none is on the Hermitian diagonal because \(k\) is odd.

Differentiating the recurrence proves, term by term,
\[
\mathcal W_{d+1}=p_d^2+a_d\mathcal W_d,
\qquad a_d=d(d-\tfrac12)=\gamma_d/\gamma_{d-1}.
\]
Indeed substitute \(p_{d+1}=xp_d-a_dp_{d-1}\) and its derivative
into \(p_{d+1}'p_d-p_d'p_{d+1}\); the two \(xp_d'p_d\) terms
cancel. Hence every resonant entry in DS5 obeys
\begin{equation}\label{dsroot:eq:resonance-update}
\frac{\mathcal W_{d+1}(u)}{\gamma_d}
-\frac{\mathcal W_d(u)}{\gamma_{d-1}}
=\frac{p_d(u)^2}{\gamma_d}.
\tag{DS6}
\end{equation}
The square here is \(p_d(u)^2\), not \(|p_d(u)|^2\): the column
node is the reflected node, whose conjugate equals the row node.
It agrees exactly with \(f_df_d^*/\gamma_d\). Off the resonance set,
the difference of the two CD quotients is the same rank-one entry by
the same polynomial telescope. Thus DS5 is also the complete CD
boundary update, including all removable denominators.

\section{Determinant and inverse updates with the original frame factors}

Introduce the nonnegative, dimensionless scalar quantities
\begin{equation}\label{dsroot:eq:leverage}
\alpha_d=\frac{g_d^*A_d^{-1}g_d}{\gamma_d},\qquad
\beta_d=\frac{f_d^*C_d^{-1}f_d}{\gamma_d}.
\tag{DS7}
\end{equation}
For any invertible matrix \(B\), the finite column expansion gives
\(\det(B+xy^*)=\det B+y^*\operatorname{adj}(B)x\).
Terms choosing the added column twice vanish because those columns are
proportional. Since \(\operatorname{adj}(B)=\det(B)B^{-1}\), this
also proves the usual rank-one determinant identity. Applying it to
DS5 gives
\begin{equation}\label{dsroot:eq:detupdate}
\boxed{\frac{\det A_{d+1}}{\det A_d}=1+\alpha_d,
\qquad\frac{\det C_{d+1}}{\det C_d}=1+\beta_d.}
\tag{DS8}
\end{equation}
Equivalently, the exact increment of \(\det A_d\) is
\(g_d^*\operatorname{adj}(A_d)g_d/\gamma_d\), with no inverse in
that numerator. A compressed determinant step is zero exactly when
\(g_d=0\); no individual nonzero step is assumed.

The complementary frame identity, with every constant present, is
\begin{equation}\label{dsroot:eq:framefactor}
\det G_d=\kappa_{\pi,Z}\frac{\det A_d}{\det C_d},
\qquad\kappa_{\pi,Z}=\frac{\det M_U}{\det W}>0.
\tag{DS9}
\end{equation}
To prove it, use \(U_0=U M_U^{-1/2}\), \(V_0=V W^{-1/2}\) only
as an auxiliary unitary frame \([U_0,V_0]\). In that frame write
\(C_d=\left(\begin{smallmatrix}B&E\\E^*&D\end{smallmatrix}\right)\).
Block elimination gives \(\det C_d=\det D\det(B-ED^{-1}E^*)\),
and direct multiplication gives the upper inverse block
\((B-ED^{-1}E^*)^{-1}\). Taking determinants and undoing the two
frame square roots proves DS9. No source metric or actual coefficient
frame has been replaced.

Set \(b_d=U^*C_d^{-1}f_d\). This is the classical rank-one inverse
update associated with Sherman and Morrison \cite{dsroot:SM}; it is not a new
inverse-matrix formula. Multiplying out the proposed inverse
checks the complete inverse update
\begin{equation}\label{dsroot:eq:inverseupdate}
C_{d+1}^{-1}=C_d^{-1}
-\frac{C_d^{-1}f_df_d^*C_d^{-1}}
{\gamma_d+f_d^*C_d^{-1}f_d},\qquad
G_{d+1}=G_d-\frac{b_db_d^*}{\gamma_d(1+\beta_d)}.
\tag{DS10}
\end{equation}
The denominator is positive. Therefore this is the exact original
kernel contraction, not an inverse of a restricted covariance.

The difference of the two scalar quantities in DS7 has a literal
complementary expression:
\begin{equation}\label{dsroot:eq:gap}
\boxed{\tau_d:=\beta_d-\alpha_d
=\frac{b_d^*G_d^{-1}b_d}{\gamma_d}\geq0.}
\tag{DS11}
\end{equation}
For a proof, the two matrices \(C_d^{1/2}V\) and \(C_d^{-1/2}U\)
have orthogonal column images, because their cross Gram is \(V^*U=0\).
Their full ranks sum to \(q\). Their two orthogonal projections sum
to the identity, which gives, after multiplying by \(C_d^{-1/2}\)
on both sides,
\[
C_d^{-1}=VA_d^{-1}V^*
+C_d^{-1}UG_d^{-1}U^*C_d^{-1}.
\]
Sandwiching by \(f_d\) and dividing by \(\gamma_d\) proves DS11.
Consequently both determinant routes, DS8--9 and DS10, agree:
\begin{equation}\label{dsroot:eq:kernelstep}
\boxed{\frac{\det G_{d+1}}{\det G_d}
=\frac{1+\alpha_d}{1+\beta_d}
=1-\frac{\tau_d}{1+\beta_d}\in(0,1].}
\tag{DS12}
\end{equation}
The equality case is exactly \(b_d=0\). Every factor
\(\kappa_{\pi,Z}\) has been retained in DS9 before its cancellation
between consecutive degrees.

The full determinant step is also the exact original relation-norm
ratio already identified in the preceding mixed derivation and the
Toda source/boundary quotient \cite{dsroot:Mixed,dsroot:Toda}%
% BEGIN VERIFIED PUBLIC PROOF LINK
~\href{https://github.com/KokunoYumeto/zeta-function-research-reader/blob/5992ad50c0f2a06921f1fcaded7b4e38097c0739/workbenches/splitzero-tandem/continuations/20260920-original-kernel-mixed-determinants/MIXED_CONFLUENT_PROOF.tex\#L1}{[proof]} \href{https://github.com/KokunoYumeto/zeta-function-research-reader/blob/5992ad50c0f2a06921f1fcaded7b4e38097c0739/workbenches/splitzero-tandem/continuations/20260920-original-kernel-mixed-determinants/providers/TODA_VOLUME_CONTROL.tex\#L1}{[proof]}%
% END VERIFIED PUBLIC PROOF LINK
{}. In physical coordinates
let \(B_n(\chi,\chi)\) be the determinant of
\(\langle S^i\chi,S^j\chi\rangle_\sigma\), let
\(D_d=\prod_{j<d}\gamma_j\), and put \(n=d-q\).
The finite mixed formula gives
\(\det C_d=|V_\chi|^2B_n/D_d\), retaining the original root
Vandermonde. Dividing at consecutive degrees yields
\[
1+\beta_d=\frac{B_{n+1}/B_n}{\gamma_d}.
\]
Gram--Schmidt for the actual multiplication source
\(P\mapsto\chi P\) gives \(B_{n+1}/B_n\) as its degree-\(n\)
monic squared norm, since the monic basis change has determinant one.
This re-expression credits the pre-existing source/boundary determinant
structure and does not introduce a different measure for the quotient.

\section{The exact four-cutoff determinant return}

For any sequence \(L_d\) define its original four-cutoff combination
\(\mathcal RL=L_q+L_{q+1}-L_{2q}-L_{2q+1}\).
The physical cutoffs are \(N=q-1,q,2q-1,2q\); hence these are the
correct kernel indices. Put
\[
w_d=\begin{cases}1&d=q\text{ or }d=2q,\\2&q<d<2q.\end{cases}
\]
Finite telescoping, first from \(q\) to \(2q\), then from \(q+1\)
to \(2q+1\), gives the exact identities
\begin{equation}\label{dsroot:eq:return}
\boxed{\begin{aligned}
\mathcal R\log\det A_d
&=-\sum_{d=q}^{2q}w_d\log(1+\alpha_d),\\
\mathcal R\log\det C_d
&=-\sum_{d=q}^{2q}w_d\log(1+\beta_d),\\
\mathcal R\log\det G_d
&=\sum_{d=q}^{2q}w_d
 \log\frac{1+\beta_d}{1+\alpha_d}
=-\sum_{d=q}^{2q}w_d
 \log\left(1-\frac{\tau_d}{1+\beta_d}\right).
\end{aligned}}
\tag{DS13}
\end{equation}
All logarithms are of positive real numbers. In product form the last
identity is
\[
\frac{\det G_q\det G_{q+1}}{\det G_{2q}\det G_{2q+1}}
=\prod_{d=q}^{2q}\left(\frac{1+\beta_d}{1+\alpha_d}\right)^{w_d}.
\]
The common frame constant from DS9 has exponent \(1+1-1-1=0\).
These formulas preserve every degree step; they do not replace the
interior steps by either endpoint alone.

\section{The actual quartet prevents complete block cancellation}

The original transported relation is
\begin{equation}\label{dsroot:eq:positive-relation}
\widetilde\chi(y)=i^{-q}\chi(c+iy)
=\prod_{a=0}^{(k-1)/2}\prod_{b=0}^{k}
\bigl((y-(2b-k)\gamma)^2+(2a-k)^2\delta^2\bigr)>0
\quad(y\in\dsrootR).
\tag{DS14}
\end{equation}
The paired roots are \(u_{ab},\overline{u_{ab}}=u_{k-a,b}\).
Each displayed factor is positive because \(2a-k\ne0\) and
\(\delta>0\). There are \(q/2\) such monic quadratic factors.
The original phase is retained: \(q=(k+1)^2\) is divisible by four,
so \(i^q=1\) on this domain. Formula DS14 is an identity for the
original relation, not a replacement relation or source norm.

For every integer \(s\geq0\), the \(q\) columns
\begin{equation}\label{dsroot:eq:tailbasis}
f_s,f_{s+1},\ldots,f_{s+q-1}
\quad\text{form a basis of }\dsrootC^q.
\tag{DS15}
\end{equation}
Here is a full proof using the original orthogonality. If a combination
\(P=\sum_{j=s}^{s+q-1}a_jp_j\) vanishes at all roots, division by
\(\widetilde\chi\) gives \(P=\widetilde\chi R\) with
\(\deg R\leq s-1\). At \(s=0\), that forces \(R=0\) directly.
For \(s>0\), orthogonality to every polynomial of degree below \(s\)
gives
\[
0=\int\overline{R(y)}P(y)\,d\sigma(y)
=\int\widetilde\chi(y)|R(y)|^2\,d\sigma(y).
\]
The last integral is strictly positive for a nonzero polynomial
\(R\), because DS14 is positive and \(d\sigma\) has positive density
on the real axis. Therefore \(R=0\), then \(P=0\), and finally all
\(a_j=0\), since the \(p_j\) have different monic leading degrees.
This proves linear independence and hence the basis claim. The
positive factor inside this one orthogonality argument is not a new
choice of quotient metric.

It follows for every \(d\geq q\) that
\begin{equation}\label{dsroot:eq:strictblocks}
C_{d+q}-C_d=\sum_{j=d}^{d+q-1}\frac{f_jf_j^*}{\gamma_j}>0,
\quad A_{d+q}-A_d>0,
\quad G_d-G_{d+q}>0.
\tag{DS16}
\end{equation}
The first strict inequality follows by evaluating its quadratic form:
a vector annihilating every squared term annihilates a basis. The second
is its congruence by the full-column-rank \(V\). For the third,
\(C_d^{-1/2}C_{d+q}C_d^{-1/2}>I\), so all eigenvalues of its inverse
are below one; congruence gives \(C_d^{-1}-C_{d+q}^{-1}>0\), and
restriction by \(U\) proves the claim.

In particular the two compressed returns in DS13 are strictly negative,
and the original kernel return is strictly positive:
\begin{equation}\label{dsroot:eq:strictreturn}
\boxed{\mathcal R\log\det A_d<0,\quad
\mathcal R\log\det C_d<0,\quad
\mathcal R\log\det G_d>0.}
\tag{DS17}
\end{equation}
For example \(\det G_q>\det G_{2q}\) and
\(\det G_{q+1}>\det G_{2q+1}\) follow from the strict positive
matrix differences in DS16, proving the last inequality directly.
These are finite strict inequalities, with no uniform lower size
assigned to any term.

There are also exact rank consequences for zero individual steps.
The columns \(g_d,\ldots,g_{d+q-1}\) span \(\dsrootC^r\), by applying
the surjection \(V^*\) to the basis DS15. At least \(r\) of these
\(q\) columns are nonzero. Summing DS10 gives
\[
G_d-G_{d+q}=\sum_{j=d}^{d+q-1}
\frac{b_jb_j^*}{\gamma_j(1+\beta_j)}>0,
\]
so \(b_d,\ldots,b_{d+q-1}\) span \(\dsrootC^m\), and at least \(m\)
of those kernel contractions are nonzero. Thus the actual invariant
image cannot force the entire original four-return to vanish, even
though an individual step may vanish according to its exact criterion.

For the complete covariance, each individual \(\beta_d\) is strictly
positive on this original domain. Here is a proof without importing a
zero-location hypothesis. The determinant of \(yI-J_d\), where
\(J_d\) is the real symmetric tridiagonal matrix with diagonal zero
and successive off-diagonals \(\sqrt{a_1},\ldots,\sqrt{a_{d-1}}\),
satisfies the same determinant expansion recurrence and initial values
as \(p_d\); hence it is \(p_d\). Every eigenvalue of \(J_d\) is
real: an eigenvector \(z\ne0\) gives
\(\lambda=z^*J_dz/(z^*z)\in\dsrootR\). Thus \(p_d\) has no nonreal
zero. Each \(u_{ab}\) has nonzero imaginary part because \(k\) is
odd, so every entry of \(f_d\) is nonzero. Positivity of
\(C_d^{-1}\) then gives \(\beta_d>0\). This does not imply a
nonzero \(g_d\) or \(b_d\) at every individual degree, because those
involve the exact original frame contractions.

\section{The complete complement correction evolves with the step}

Keep the fixed original matrices
\[
Y=\operatorname{diag}(u_{ab}),\quad
P=VW^{-1}V^*,\quad Q=I-P=UM_U^{-1}U^*,\quad
H=(V^*YV)W^{-1},\quad F=V^*YU,
\]
and define
\[
T_d=U^*C_dV,\quad e_d=U^*f_d,\quad
\rho_d=FM_U^{-1}e_d=V^*YQf_d,
\quad \mathcal E_d=FM_U^{-1}T_d-T_d^*M_U^{-1}F^*.
\]
The full CD displacement, proved by the recurrence telescope, is
\(YC_d-C_dY^*=(f_df_{d-1}^*-f_{d-1}f_d^*)/\gamma_{d-1}\).
Inserting \(I=P+Q\) and multiplying by \(V^*,V\) gives the exact
compressed identity
\begin{equation}\label{dsroot:eq:correctedCD}
HA_d-A_dH^*
=\frac{g_dg_{d-1}^*-g_{d-1}g_d^*}{\gamma_{d-1}}
-\mathcal E_d.
\tag{DS18}
\end{equation}
Indeed \(HA_d=V^*YPC_dV\) and
\(A_dH^*=V^*C_dPY^*V\); the discarded complement difference would
be precisely \(\mathcal E_d\). It is not discarded here.
Applying DS5 to \(T_d\) proves
\begin{equation}\label{dsroot:eq:correctionstep}
\boxed{T_{d+1}=T_d+\frac{e_dg_d^*}{\gamma_d},\qquad
\mathcal E_{d+1}-\mathcal E_d
=\frac{\rho_dg_d^*-g_d\rho_d^*}{\gamma_d}.}
\tag{DS19}
\end{equation}
All cross terms and their signs are present.

Projection of the original recurrence gives
\(g_{d+1}=Hg_d-a_dg_{d-1}+\rho_d\).
Writing this right side as \(\xi_d\), the exact next scalar update can
also be obtained using only the current inverse and the full correction
data:
\begin{equation}\label{dsroot:eq:nextalpha}
\alpha_{d+1}=\frac1{\gamma_{d+1}}
\left(\xi_d^*A_d^{-1}\xi_d
-\frac{|g_d^*A_d^{-1}\xi_d|^2}
{\gamma_d+g_d^*A_d^{-1}g_d}\right),
\qquad\xi_d=Hg_d-a_dg_{d-1}+\rho_d.
\tag{DS20}
\end{equation}
This follows by applying the directly checked inverse formula DS10
to \(A_{d+1}\). Omitting \(\rho_d\) would be an additional,
unproved cancellation; DS20 instead retains it exactly.

\section{The literal source-scaling map in the original coefficient algebra}

We now calculate what the invariant-image identity actually cancels.
Retain the OCF period matrix \(\mathsf H\), the original ordered
variables \((x_1,x_2,x_3,x_4)\), and the Segre substitution
\[
t=(z_0w_0,z_0w_1,z_1w_0,z_1w_1)^T,\quad x=\mathsf Ht,
\quad Q(x)=Q_0(\mathsf H^{-1}x),\quad
Q_0(t)=t_1t_4-t_2t_3.
\]
For a homogeneous polynomial \(F\) of degree \(s\), let
\(\psi_sF=F(\mathsf Ht)\), a biform of bidegree \((s,s)\).
OCF3 proves that its exact kernel is \(Q\) times degree \(s-2\)
polynomials, and that every biform has a polynomial lift. To see the
surjection directly, the monomial
\(z_0^a z_1^{s-a}w_0^b w_1^{s-b}\) is the image of
\(t_1^h t_2^{a-h}t_3^{b-h}t_4^{s-a-b+h}\), where
\(\max(0,a+b-s)\leq h\leq\min(a,b)\).

Define the actual degree-preserving biform derivation and its
period-coordinate matrix by
\begin{equation}\label{dsroot:eq:generator}
\mathcal Y=\gamma(w_0\partial_{w_0}-w_1\partial_{w_1})
-i\delta(z_0\partial_{z_0}-z_1\partial_{z_1}),\quad
\Lambda=\operatorname{diag}(\gamma-i\delta,-\gamma-i\delta,
\gamma+i\delta,-\gamma+i\delta),\quad
\mathsf L=\mathsf H\Lambda\mathsf H^{-1}.
\tag{DS21}
\end{equation}
For any matrix \(B\), define the polynomial derivation
\(\dsrootcalD_BF=(Bx)\cdot\nabla F\). The chain rule proves the exact
commuting identity
\begin{equation}\label{dsroot:eq:generator-map}
\psi_s\dsrootcalD_{\mathsf L}=\mathcal Y\psi_s.
\tag{DS22}
\end{equation}
Indeed \(\mathcal Yt=\Lambda t\), so
\(\mathcal Y(\mathsf Ht)=\mathsf H\Lambda t=\mathsf Lx\).
The two sums \(\Lambda_{11}+\Lambda_{44}\) and
\(\Lambda_{22}+\Lambda_{33}\) are both zero, so
\(\dsrootcalD_{\mathsf L}Q=0\); the derivation is well defined modulo
the original \(Q\). On the ordered biform coefficient basis of degree
\(k\), \(\mathcal Y\) has diagonal matrix exactly \(Y\) from DS1.
This proves the coordinate-level morphism between the physical source
scaling and the coefficient generator; no new action is assumed.

\section{The invariant average cancels the diagonal generator exactly}

Let \(D=\operatorname{diag}(\zeta,\zeta^2,\zeta^3,\zeta^4)\),
\(\zeta=e^{2\pi i/5}\), let \(gF(x)=F(D^{-1}x)\), and retain
the exact original averaging operator
\(\mathsf R=\tfrac15\sum_{a=0}^4g^a\).
The original invariant space is \(\mathscr I_k=\{F:gF=F\}\), and
the original image is \(J_k=\psi_k\mathscr I_k=\ker\pi\).
By the chain rule,
\[
g^a\dsrootcalD_Bg^{-a}=\dsrootcalD_{D^aBD^{-a}}.
\]
For invariant \(F\), averaging therefore gives
\begin{equation}\label{dsroot:eq:reynolds}
\mathsf R\dsrootcalD_{\mathsf L}F=\dsrootcalD_{\mathsf L_0}F,
\qquad
\mathsf L_0=\frac15\sum_{a=0}^4D^a\mathsf LD^{-a}
=\operatorname{diag}(\mathsf L_{11},\ldots,\mathsf L_{44}).
\tag{DS23}
\end{equation}
For the final equality the \((i,j)\) entry is multiplied by
\(\tfrac15\sum_a\zeta^{a(i-j)}\), which is one for \(i=j\)
and zero for the distinct weights \(i\ne j\). Thus the average is
not being estimated: it is computed exactly. Its value is invariant,
so its image under \(\psi_k\) is annihilated by \(\pi\).
Writing \(\mathsf L_{\mathrm{off}}=\mathsf L-\mathsf L_0\), we obtain
the exact original-image cancellation
\begin{equation}\label{dsroot:eq:offdiag-cancellation}
\boxed{\pi\psi_k\dsrootcalD_{\mathsf L}F
=\pi\psi_k\dsrootcalD_{\mathsf L_{\mathrm{off}}}F
\quad(F\in\mathscr I_k).}
\tag{DS24}
\end{equation}
Every actual diagonal entry is present before this cancellation, and
every actual off-diagonal entry is retained afterwards.

Here are complete invariant certificates for the exact columns of
\(Z\), so DS24 is an explicit coefficient calculation, not a
statement about an unspecified invariant lift. Put
\(\mathcal A=\prod_{a=1}^4g^aQ\) and
\(A=\psi_8\mathcal A\). For any biform \(h\) of degree \(k-8\),
choose its original monomial lift \(\widetilde h\) through the explicit
Segre monomial formula in the preceding section. Then
\[
F_h=5\mathsf R(\mathcal A\widetilde h)\in\mathscr I_k,
\qquad\psi_kF_h=A h.
\]
Indeed, modulo \(Q\), the four terms \(g^a\mathcal A\) for
\(a\ne0\) each contain \(Q\); the term \(a=0\) is precisely
\(\mathcal A\). This proves the original factor 5 and supplies the
certificates for \(M_A\). For a retained column of \(Z_k\), OCF20
supplies an invariant \(F_j\) with
\(N_k\psi_kF_j=(Z_k)_j\). OCF7 gives
\(\psi_kF_j=S_k(Z_k)_j+A h_j\), with its explicitly recorded
division coefficients \(h_j\). Subtracting
\(5\mathsf R(\mathcal A\widetilde h_j)\) from \(F_j\) supplies
an invariant certificate for the exact column \(S_k(Z_k)_j\).
No chosen frame coefficient is altered.

Let these complete certificates, in the original column order, be
\(\mathcal F_1,\ldots,\mathcal F_r\). The surviving coupling is the
fully specified finite matrix
\begin{equation}\label{dsroot:eq:omega}
\boxed{\Omega:=\pi YZ
=\left[\pi\psi_k
 \sum_{i\ne j}\mathsf L_{ij}x_j\partial_{x_i}\mathcal F_\ell
 \right]_{\ell=1}^r,
\qquad F=\Omega^T.}
\tag{DS25}
\end{equation}
The last equality is \(V^*YU=Z^TY\pi^T=(\pi YZ)^T\), using
the diagonal matrix \(Y^T=Y\), not a conjugation. Formula DS25 is
an exact expression in the original period entries, original invariant
certificates and the already specified elimination maps. It reduces
the source coupling to the off-diagonal generator, without claiming
that generator or its quotient image vanishes.

The conductor part has a further exact product-rule cancellation:
\begin{equation}\label{dsroot:eq:conductor-cancel}
\pi Y M_A=\pi M_{\mathcal Y A}.
\tag{DS26}
\end{equation}
Indeed \(\mathcal Y(Ah)=(\mathcal YA)h+A\mathcal Yh\), and the
second term lies in \(\operatorname{im}M_A\subseteq\ker\pi\).
The first term in DS26 is retained. Substituting DS25 into DS19 gives
the complete original correction and its step increment solely with
the surviving coefficient matrix:
\begin{equation}\label{dsroot:eq:omega-correction}
\mathcal E_d=\Omega^TM_U^{-1}T_d-T_d^*M_U^{-1}\overline\Omega,
\qquad \rho_d=\Omega^TM_U^{-1}e_d.
\tag{DS27}
\end{equation}
These formulas determine the correction at the actual period. They do
not replace it by zero. Its exact vanishing criterion is the equality
of the two displayed products in DS27.

\section{The precise scope of the invariant cancellation}

The following coefficient test proves that cyclic invariance and the
five-orbit condition alone do not erase the off-diagonal term in DS24.
It is not assigned to the actual period, used as its replacement frame,
or used to infer a determinant asymptotic. It tests the claimed
algebraic implication itself.

Take the invertible coefficient matrix
\(\mathsf H_*=I+E_{12}\), so
\(x_1=t_1+t_2,x_2=t_2,x_3=t_3,x_4=t_4\). Then
\[
Q_*=(x_1-x_2)x_4-x_2x_3,
\qquad g^aQ_*=x_1x_4-x_2x_3-\zeta^{-a}x_2x_4.
\]
The five quadrics are pairwise nonassociate: the coefficient of
\(x_1x_4\) fixes any associate scalar at one, and their
\(x_2x_4\) coefficients differ. They are prime rank-four quadrics
because they are invertible linear transforms of the Segre equation.
Thus the same five-orbit coefficient construction applies.
The transported generator has
\(\mathsf L_{*,\mathrm{off}}=-2\gamma E_{12}\).
At the permitted degree \(k=9\), the monomial
\(F_*=x_1^8x_2\) is invariant because \(8+2\equiv0\pmod5\), but
\[
\dsrootcalD_{\mathsf L_{*,\mathrm{off}}}F_*
=-16\gamma x_1^7x_2^2.
\]
Its restriction is not in the invariant image. To prove this, set
\(z_1=0\) after the Segre substitution. Then \(x_3=x_4=0\), and
the only degree-nine invariant monomials which survive are
\(x_1^8x_2\) and \(x_1^3x_2^6\): with exponent \(b\) on
\(x_2\), the weight is \(9+b\), so \(b=1\) or \(6\).
Set \(z_0=w_1=1\) and write \(s=w_0+1\). Their two restrictions
are \(s^8,s^3\), whereas the derivative restriction is a nonzero
multiple of \(s^7\). Distinct monomial degrees make the latter
independent of the former two. Thus its quotient image is nonzero.
This proves that the off-diagonal remainder cannot be deleted as a
formal consequence of the averaging identity. At the actual period
its value remains the exact matrix DS25, not this test's value.

Independently of whether the actual correction has further vanishing
entries, DS15--17 apply to the actual quartet and original full-rank
frame. They prove that the complete four-cutoff kernel return cannot
cancel to zero. They give no lower magnitude at the \(kq\) scale.

\section{A complete nonzero-correction degree-step example}

This second example tests the determinant and correction formulas with
the exact Gamma mass; its two-node coefficient frame is explicitly an
algebra fixture, not an arithmetic-period substitute. Take
\[
u=(0,i),\quad V=(1,i)^T,\quad Z=(1,-i)^T,
\quad \pi=(i,1),\quad U=(i,1)^T,
\quad M=M_\sigma.
\]
The matrices at degrees \(d=2,3\) are
\[
C_2=M^{-1}\begin{pmatrix}1&1\\1&3\end{pmatrix},\qquad
C_3=M^{-1}\begin{pmatrix}7/6&3/2\\3/2&9/2\end{pmatrix}.
\]
Here \(f_2=(-1/2,-3/2)^T\), \(\gamma_2=3M/2\), and direct
multiplication verifies \(C_3-C_2=f_2f_2^*/\gamma_2\). The original
frame factors in this fixture are \(W=M_U=2\), so
\(\kappa_{\pi,Z}=1\). The complete scalars are
\begin{equation}\label{dsroot:eq:fixture}
\begin{gathered}
A_2=4/M,\quad A_3=17/(3M),\quad
\det C_2=2/M^2,\quad\det C_3=3/M^2,\\
G_2=2M,\quad G_3=17M/9,\quad
\alpha_2=5/12,\quad\beta_2=1/2,\quad\tau_2=1/12,\\
\frac{A_3}{A_2}=17/12,\quad
\frac{\det C_3}{\det C_2}=3/2,\quad
\frac{G_3}{G_2}=17/18.
\end{gathered}\tag{DS28}
\end{equation}
For the same fixture \(H=i/2\). The exact corrections are
\(\mathcal E_2=2i/M\) and \(\mathcal E_3=10i/(3M)\), giving
the nonzero increment \(4i/(3M)\). These follow either by direct
use of DS18 or from DS19 with the exact \(f_2,g_2,e_2\).
Thus the determinant update is not relying on a vanished correction.
The test fixture does not satisfy the quartet's conjugate-pair relation
DS14, and it is not used to prove the strict block theorem; that theorem
was proved directly for the actual quartet.

\section{Mass, physical phases, and final boundary of the calculation}

Under the diagnostic scaling \(d\sigma\mapsto a\,d\sigma\),
\(a>0\), the monic polynomials, frames and root coordinates stay fixed,
while \(\gamma_j\mapsto a\gamma_j\),
\(C_d,A_d,T_d,\mathcal E_d\mapsto a^{-1}\) times themselves,
\(G_d\mapsto aG_d\), and \(b_d\mapsto ab_d\).
Every \(\alpha_d,\beta_d,\tau_d\) and every one-degree determinant
ratio is unchanged. DS9 scales by \(a^{q-r}=a^m\), exactly the
determinant scale of the original rank-\(m\) restricted metric. Its
constant mass factor cancels in DS13 only after being retained at each
cutoff. The actual source has \(a=1\), \(M_\sigma=\sqrt{2\pi}\).

The physical \(S\)-evaluation columns are \(i^jf_j\), and their
compressed columns are \(i^jg_j\). The phases multiply to one in every
rank-one covariance update and every scalar in DS7--12. The physical
coefficient generator on degree \(k\) is
\(cI+iY\). Its coupling is
\(\pi(cI+iY)Z=i\Omega\), because the complete original identity
\(\pi Z=0\) removes the scalar term. In polynomial coordinates its
derivation is \(\dsrootcalD_{\frac12I+i\mathsf L}\); the radial part acts
as \(k/2=c\) at degree \(k\). These are exact coordinate maps,
not changes of the original scaling operator.

The finite outcomes are DS8, DS10--13, the strict block and return
statements DS15--17, and the fully specified surviving invariant-image
coupling DS24--27. The original period has not been numerically sampled.
No source-index, rendering, packaging or publication action is part of
this derivation, and no estimate of the unresolved coefficient is
deduced from an algebra fixture.

\begin{thebibliography}{9}
\bibitem{dsroot:SM} J. Sherman and W. J. Morrison,
\emph{Adjustment of an Inverse Matrix Corresponding to a Change in One
Element of a Given Matrix}, Annals of Mathematical Statistics
21 (1950), 124--127,
\url{https://doi.org/10.1214/aoms/1177729893}.
Historical attribution; bibliographic metadata verified, original article
not inspected. The complete rank-one calculation used here is given at DS10.
\bibitem{dsroot:AV} G. Akemann and G. Vernizzi,
\emph{Characteristic Polynomials of Complex Random Matrix Models},
arXiv:hep-th/0212051v2, original author TeX, \texttt{CD} and its
algebraic-continuation footnote,
\url{https://arxiv.org/abs/hep-th/0212051v2}.
\bibitem{dsroot:DLMF} T. H. Koornwinder, R. Wong, R. Koekoek and
R. F. Swarttouw; W. P. Reinhardt, NIST DLMF Chapter 18,
original retained TeX:
\url{https://dlmf.nist.gov/18.19.E8},
\url{https://dlmf.nist.gov/18.19.E9a},
\url{https://dlmf.nist.gov/18.23.E7}.
\bibitem{dsroot:OCF} Split-Zero programme,
\emph{An exact conductor frame and a fixed invariant-generator
recursion}, original OCF1--26, especially OCF3, OCF7, OCF17,
and OCF20--24. All period and amplitude coefficients are retained.

% BEGIN VERIFIED PUBLIC PROOF LINK
\href{https://github.com/KokunoYumeto/zeta-function-research-reader/blob/5992ad50c0f2a06921f1fcaded7b4e38097c0739/workbenches/splitzero-tandem/continuations/20260920-original-kernel-mixed-determinants/providers/ORIGINAL_CONDUCTOR_STRIP_FRAME.tex\#L27}{Source proof OCF1}.
% END VERIFIED PUBLIC PROOF LINK
\bibitem{dsroot:MR} Split-Zero programme,
\emph{The original observation kernel as an exact mixed-relation
determinant}, original MR1--23, especially MR8--11 and MR22--23.

% BEGIN VERIFIED PUBLIC PROOF LINK
\href{https://github.com/KokunoYumeto/zeta-function-research-reader/blob/5992ad50c0f2a06921f1fcaded7b4e38097c0739/workbenches/splitzero-tandem/continuations/20260920-original-kernel-mixed-determinants/ORIGINAL_KERNEL_MIXED_DETERMINANT.tex\#L36}{Source proof MR1}.
% END VERIFIED PUBLIC PROOF LINK
\bibitem{dsroot:CD} Split-Zero programme,
\emph{Two boundary polynomials for the original Gamma kernel and the
exact correction after coefficient compression}, complete CD1--22,
20 September 2026, frozen preceding local derivation.

% BEGIN VERIFIED PUBLIC PROOF LINK
\href{https://github.com/KokunoYumeto/zeta-function-research-reader/blob/main/workbenches/splitzero-tandem/continuations/20260921-kernel-degree-steps/COMPLETE_CD_BOUNDARY.tex\#L42}{Source proof CD1}.
% END VERIFIED PUBLIC PROOF LINK
\bibitem{dsroot:Mixed} Split-Zero programme,
\emph{The finite mixed and confluent Christoffel determinant on the
original cyclic source}, complete preceding finite proof, especially
its main mixed identity and original relation-norm ratio.

% BEGIN VERIFIED PUBLIC PROOF LINK
\href{https://github.com/KokunoYumeto/zeta-function-research-reader/blob/5992ad50c0f2a06921f1fcaded7b4e38097c0739/workbenches/splitzero-tandem/continuations/20260920-original-kernel-mixed-determinants/MIXED_CONFLUENT_PROOF.tex\#L1}{Source proof}.
% END VERIFIED PUBLIC PROOF LINK
\bibitem{dsroot:Toda} User-supplied programme continuation,
\emph{Source--boundary volume dynamics for the arithmetic control},
12 September 2026, original \texttt{NOTE.tex}, equations (11)--(16)
and (19). The earlier source/boundary determinant quotient remains
its original result.

% BEGIN VERIFIED PUBLIC PROOF LINK
\href{https://github.com/KokunoYumeto/zeta-function-research-reader/blob/5992ad50c0f2a06921f1fcaded7b4e38097c0739/workbenches/splitzero-tandem/continuations/20260920-original-kernel-mixed-determinants/providers/TODA_VOLUME_CONTROL.tex\#L1}{Source proof}.
% END VERIFIED PUBLIC PROOF LINK
\end{thebibliography}

\endgroup

% END COMPLETE CD AND DEGREE STEP 20260921
